% Part:sets-functions-relations
% Chapter: size-of-sets
% Section: pairing-alt

\documentclass[../../../include/open-logic-section]{subfiles}

\begin{document}

\olfileid{sfr}{siz}{pai-alt}

\olsection{మరొక జతీకరణ ప్రమేయం}

\begin{explain}
విలోమాలను సులభంగా కనుక్కోగలిగేలా $\Nat^2$ను లెక్కించే ఇతర
పద్ధతులూ ఉన్నాయి. ఇదొకటి. లెక్కింపును పట్టికగా ఊహించడానికి బదులు,
(మొదట ఖాళీగా ఉన్న) స్థానాలతో జతచేసిన ధన పూర్ణ సంఖ్యల జాబితాతో
మొదలుపెట్టండి. ఆ స్థానాలను $\tuple{n,m}$ జతలతో వరుసగా ఇలా నింపుతున్నట్లు
ఊహించండి. మొదటి స్థానంలో $0$ ఉన్న జతలతో (అంటే $\tuple{0,m}$ జతలతో)
మొదలుపెట్టి, మొదటిదైన $\tuple{0,0}$ను మొదటి ఖాళీ స్థానంలో పెట్టండి;
తర్వాత ఒక ఖాళీ స్థానాన్ని వదిలి, రెండవదైన $\tuple{0,1}$ను తదుపరి
ఖాళీ స్థానంలో పెట్టండి; మళ్లీ ఒకటి వదిలి ఇదే విధంగా కొనసాగండి.
ఇప్పుడు మన అసంపూర్ణ లెక్కింపు ఆరంభం ఇలా ఉంటుంది:
\[\small
\begin{array}{@{}c c c c c c c c c c c@{}}
\mathbf 1 & \mathbf 2 & \mathbf 3 & \mathbf 4 & \mathbf 5 & \mathbf 6 & \mathbf 7 & \mathbf 8 & \mathbf 9 & \mathbf{10} & \dots \\ \\
\tuple{0,0} &  & \tuple{0,1} &  & \tuple{0,2} &  & \tuple{0,3} & & \tuple{0,4} &  & \dots \\
\end{array}
\]
ఇంకా ఖాళీగా ఉన్న స్థానాల్లో $\tuple{1,m}$ జతలతో ఇదే ప్రక్రియను
చేయండి; ప్రతి రెండవ ఖాళీ స్థానాన్ని మళ్లీ వదిలివేయండి:
\[\small
\begin{array}{@{}c c c c c c c c c c c@{}}
\mathbf 1 & \mathbf 2 & \mathbf 3 & \mathbf 4 & \mathbf 5 & \mathbf 6 & \mathbf 7 & \mathbf 8 & \mathbf 9 & \mathbf{10} & \dots \\ \\
\tuple{0,0} & \tuple{1,0} & \tuple{0,1} &  & \tuple{0,2} & \tuple{1,1} &
\tuple{0,3} & & \tuple{0,4} &  \tuple{1,2} & \dots \\
\end{array}
\]
ఇదే విధంగా $\tuple{2,m}$, $\tuple{3,m}$ మొదలైన జతలను చేర్చండి.
పూర్తయిన మన లెక్కింపు ఇలా మొదలవుతుంది:
\[\small
\begin{array}{@{}cc c c c c c c c c c@{}}
\mathbf 1 & \mathbf 2 & \mathbf 3 & \mathbf 4 & \mathbf 5 & \mathbf 6 & \mathbf 7 & \mathbf 8 & \mathbf 9 & \mathbf{10} & \dots \\ \\
\tuple{0,0} & \tuple{1,0} & \tuple{0,1} & \tuple{2,0}  & \tuple{0,2} &
\tuple{1,1} & \tuple{0,3} & \tuple{3,0}  & \tuple{0,4} &  \tuple{1,2} & \dots \\
\end{array}
\]
పై లెక్కింపును అనుసరించి పట్టిక కణాలకు సంఖ్యలు ఇస్తే, చక్కని
జిగ్‌జాగ్ గీత కనిపించదు; బదులుగా ఈ అమరిక వస్తుంది:
\[
\begin{array}{ c | c | c | c | c | c | c | c }
& \mathbf 0 & \mathbf 1 & \mathbf 2 & \mathbf 3 & \mathbf 4 & \mathbf 5 & \dots \\
\hline
\mathbf 0 & 1 & 3 & 5 & 7 & 9 & 11 & \dots \\
\hline
\mathbf 1 & 2 & 6 & 10 & 14 & 18 & \dots & \dots \\
\hline
\mathbf 2 & 4 & 12 & 20 & 28 & \dots & \dots & \dots \\
\hline
\mathbf 3 & 8 & 24 & 40 & \dots & \dots & \dots & \dots \\
\hline
\mathbf 4 & 16 & 48 & \dots & \dots & \dots & \dots & \dots \\
\hline
\mathbf 5 & 32 & \dots & \dots & \dots & \dots & \dots & \dots \\
\hline
\vdots & \vdots & \vdots & \vdots & \vdots & \vdots & \vdots & \ddots\\
\end{array}
\]
$0$వ అడ్డువరుసలోని జతలు మన లెక్కింపులో బేసి సంఖ్య గల స్థానాల్లో
ఉన్నాయని చూడవచ్చు; అంటే $\tuple{0,m}$ జత $2m+1$వ స్థానంలో ఉంటుంది.
రెండవ అడ్డువరుసలోని $\tuple{1,m}$ జతలు ఒక బేసి సంఖ్యకు రెట్టింపు
అయిన స్థానాల్లో, ప్రత్యేకంగా $2 \cdot (2m+1)$వ స్థానాల్లో ఉంటాయి.
మూడవ అడ్డువరుసలోని $\tuple{2,m}$ జతలు ఒక బేసి సంఖ్యకు నాలుగు రెట్లు
అయిన $4 \cdot (2m+1)$వ స్థానాల్లో ఉంటాయి; ఇలాగే కొనసాగుతుంది.
ఒక్కో అడ్డువరుసలోని $(2m+1)$ గుణకాలు $1$, $2$, $4$, $8$, \dots;
ఇవి కచ్చితంగా $2$ యొక్క ఘాతాలు: $1= 2^0$, $2 = 2^1$,
$4 = 2^2$, $8 = 2^3$, \dots\@ వాస్తవానికి, సంబంధిత ఘాతం ఎప్పుడూ
ఆ జతలోని మొదటి సంఖ్యే. కాబట్టి $\tuple{n,m}$ జతకు గుణకం $2^n$.
దీంతో సాధారణ సూత్రం $2^n \cdot (2m+1)$ వస్తుంది. అయితే ఇది జతలను
\emph{ధన} పూర్ణ సంఖ్యలకు పంపుతుంది; అంటే $\tuple{0,0}$ స్థానం $1$.
స్థానం $0$తో మొదలుపెట్టాలంటే ఫలితం నుంచి $1$ తీసివేయాలి. అప్పుడు:
\sourcecorrection{OLTESIZ-004}{ఆంగ్ల మూలం రెండవ $\tuple{0,m}$ జతను
$\tuple{0,2}$ అంటుంది; పట్టిక, విధానం ప్రకారం అది $\tuple{0,1}$.
దానిని సరిచేశాం.}
\sourcecorrection{OLSIZ-003}{ఆంగ్ల మూలం వరుసగా చేర్చాల్సిన జతలను
$\tuple{2,m}$, $\tuple{2,m}$ అని పునరావృతం చేస్తుంది; అడ్డువరుసల
క్రమం ప్రకారం రెండవదాన్ని $\tuple{3,m}$గా సరిచేశాం.}
\end{explain}

\begin{ex}
$\Nat^2$ సహజ సంఖ్యల జతల సమితికి
\[
h(n,m) = 2^n (2m+1) - 1
\]
అని ఇచ్చిన $h\colon \Nat^2 \to \Nat$ ఒక జతీకరణ ప్రమేయం.
\end{ex}

\begin{explain}
అందువల్ల $\Nat^2$కు మన రెండవ లెక్కింపులో $\tuple{0,0}$ సంకేత సంఖ్య
$h(0,0) = 2^0(2\cdot 0+1) - 1 = 0$; $\tuple{1,2}$ సంకేత సంఖ్య
$2^{1} \cdot (2 \cdot 2 + 1) - 1 = 2 \cdot 5 - 1 = 9$;
$\tuple{2,6}$ సంకేత సంఖ్య $2^{2} \cdot (2 \cdot 6 + 1) - 1 = 51$.
\end{explain}

కొన్నిసార్లు సంకేతీకరణ సంగ్రస్తం కావాలని కోరకుండానే సహజ సంఖ్యల
$\Nat^2$ జతలను సంకేతీకరిస్తే సరిపోతుంది. అటువంటి సంకేతీకరణల విలోమాలు
పాక్షిక ప్రమేయాలు మాత్రమే.

\begin{ex}
\[
j(n,m) = 2^n3^m
\]
అని ఇచ్చిన $j\colon \Nat^2 \to \Nat^+$ అనేది $\Nat^2 \to \Nat$
\tetoken{ఒక ఒకటి-ఒకటి}{!!a{injective}} ప్రమేయం.
\end{ex}

\end{document}
